\section{はじめに}\label{sec:introduction}

頻出パターンマイニングはデータマイニングの中核的手法として広く研究されてきた~\cite{agrawal1994fast,zaki2000scalable}。
特に、トランザクションデータベースにおける頻出アイテムセットの発見は、
マーケットバスケット分析やウェブログ解析など多くの応用分野で活用されている。

近年、トランザクションの時間的構造に注目し、
アイテムセットのサポートが特定期間に集中する\textbf{密集区間}（dense interval）を
検出する手法が提案されている。
スライディングウィンドウを用いた密集アイテムセットマイニングは、
「いつ」パターンが頻出するかという時間的情報を提供する点で、
従来の頻出パターンマイニングを補完する。

しかしながら、既存の密集区間マイニングは本質的に\textbf{記述的}（descriptive）な手法であり、
過去のデータから密集区間を検出するにとどまる。
実応用においては、「次にいつ密集区間が発生するか」「どの程度持続するか」という
\textbf{予測的}（predictive）な問いに答えることが求められる場面が多い。
例えば、小売業における需要急増の予測、ソーシャルメディアにおけるバイラル拡散の予測、
通信ネットワークにおけるトラフィックバーストの予測などである。

本研究では、密集区間マイニングの出力を点過程・生存分析の枠組みと接続し、
記述的パターンマイニングから予測的パターンマイニングへの橋渡しを行う。
具体的には、以下の3つの概念を導入する。

\begin{enumerate}
  \item \textbf{密集区間発生過程}（Dense Interval Occurrence Process; DIOP）：
    密集区間の発生を自己励起点過程としてモデル化する枠組み
  \item \textbf{密集区間間時間}（Inter-Dense Interval Time; IDIT）：
    連続する密集区間の間隔を確率変数として捉え、その分布特性を分析する概念
  \item \textbf{密集持続時間予測}（Dense Duration Prediction）：
    生存分析に基づき、密集区間の持続時間を予測する手法
\end{enumerate}

本論文の貢献は以下の通りである。
\begin{itemize}
  \item 密集区間マイニングと時間的点過程を統合する新たな理論的枠組みDIOPの提案
  \item IDITの分布特性に関する体系的な分析手法の確立
  \item Hawkes過程に基づく次回密集区間発生時刻の予測モデル（Hawkes-Dense）の開発
  \item Weibull回帰に基づく密集区間持続時間の予測手法の提案
  \item 合成データによる包括的な実験評価
\end{itemize}
